\writeSectionFrame{\presentationTitle}{Fazit}
\begin{frame}{Anwendungsfälle}
	\begin{itemize}
 		\item Ein Objekt kann verschiedene Zustände annehmen.
 		\item Zustände unterscheiden sich voneinander
 		\item Ein Objekt kann zu einem Zeitpunkt in nur einem Zustand sein
 		\item Weitere Zustände können hinzukommen
 		\item Zustände sollen unabhängig voneinander geändert werden können
 	\end{itemize}
\end{frame}
\note{
	Um zu entscheiden, ob sich dieses Muster lohnt, muss es sich lohnen, an die Stelle
	von Fallunterscheidungen eigene Klassen zu setzen.
}

\begin{frame}{Für welche Anwendungen ist das Zustand-Muster geeignet?}
	\begin{itemize}
 		\item UI-Komponenten
 		\item Automaten
 		\item Workflow-Systeme
 		\item Steuerungssysteme
 		\item Systeme mit hohem Interaktionsgrad (z.B. Spiele)
 	\end{itemize}
\end{frame}

\vorteil{Klare Trennung der Zustandslogik}
\note{
    Zustandsabhängiges Verhalten wird in eigenen Klassen gekapselt, was den Code übersichtlich und wartbar macht. Dadurch wird vermieden, dass Zustandslogik in if- oder switch-Statements verstreut ist.
}

\vorteil{Einfaches Hinzufügen neuer Zustände}
\note{
    Neue Zustände können einfach durch das Hinzufügen neuer Klassen oder Objekte implementiert werden, ohne die bestehende Logik zu ändern (Open-Closed-Prinzip).
}

\vorteil{Zustandsübergänge explizit und modular}
\note{
    Die Übergänge zwischen Zuständen werden explizit durch die Logik der Zustandsklassen definiert. Dies macht es einfacher, die Zustandsübergänge nachzuvollziehen und zu warten.
}

\vorteil{Verbesserte Testbarkeit}
\note{
    Da die Zustände als separate Klassen implementiert sind, können diese leicht einzeln getestet werden, was die Testbarkeit des Systems erhöht.
}

\vorteil{Flexibilität bei der Zustandsverwaltung}
\note{
    Jeder Zustand kann sich selbst ändern oder den Zustand der übergeordneten Klasse modifizieren, ohne dass die Hauptklasse sich um die Details kümmern muss. Dies schafft Flexibilität bei der Zustandsverwaltung.
}

\vorteil{Erweiterbarkeit}
\note{
    Neue Zustände können durch Hinzufügen von neuen Klassen einfach integriert werden, ohne dass bestehender Code stark modifiziert werden muss.
}

\nachteil{Erhöhte Anzahl von Klassen}
\note{
    Jede Zustandsänderung erfordert eine neue Klasse, was zu einer Vielzahl von Klassen führen kann. Dies kann bei vielen Zuständen den Code komplexer und schwieriger zu pflegen machen, besonders in kleinen Projekten.
}

\nachteil{Komplexität bei einfachen Problemen}
\note{
    Für einfache Zustandswechsel kann das Zustandsmuster übertrieben sein. In solchen Fällen kann es zu unnötiger Komplexität führen, die mit if-Anweisungen leichter zu bewältigen wäre.
}

\nachteil{Zustandsübergänge können schwer nachvollziehbar sein}
\note{
    Da die Zustandsübergänge in den Zustandsklassen selbst verwaltet werden, kann es bei einer Vielzahl von Zuständen schwierig sein, den Überblick über alle möglichen Übergänge und deren Bedingungen zu behalten.
}

\nachteil{Speicher- und Performance-Overhead}
\note{
    Da für jeden Zustand eine eigene Klasse erstellt wird, kann dies zu einem zusätzlichen Speicher- und Performance-Overhead führen, insbesondere bei Systemen mit vielen Zuständen oder ressourcenbeschränkten Umgebungen.
}

\nachteil{Potenzial für überkomplizierte Strukturen}
\note{
    Wenn es viele Zustände gibt und die Zustandsübergänge komplex sind, kann das System unnötig kompliziert werden. Es kann schwierig werden, die gesamte Zustandslogik zu verstehen, wenn diese nicht gut dokumentiert ist.
}

\nachteil{Erfordert enge Zusammenarbeit der Zustandsklassen}
\note{
    Da die Zustandsklassen eng miteinander verknüpft sind, müssen sie oft über Details der anderen Zustände Bescheid wissen, was zu einer gewissen Kopplung führt. Dies kann gegen das Prinzip der losen Kopplung verstoßen.
}